% Options for packages loaded elsewhere
\PassOptionsToPackage{unicode}{hyperref}
\PassOptionsToPackage{hyphens}{url}
\PassOptionsToPackage{dvipsnames,svgnames,x11names}{xcolor}
%
\documentclass[
  11pt,
  letterpaper,
]{article}
\usepackage{amsmath,amssymb}
\usepackage{setspace}
\usepackage{iftex}
\ifPDFTeX
  \usepackage[T1]{fontenc}
  \usepackage[utf8]{inputenc}
  \usepackage{textcomp} % provide euro and other symbols
\else % if luatex or xetex
  \usepackage{unicode-math} % this also loads fontspec
  \defaultfontfeatures{Scale=MatchLowercase}
  \defaultfontfeatures[\rmfamily]{Ligatures=TeX,Scale=1}
\fi
\usepackage{lmodern}
\ifPDFTeX\else
  % xetex/luatex font selection
\fi
% Use upquote if available, for straight quotes in verbatim environments
\IfFileExists{upquote.sty}{\usepackage{upquote}}{}
\IfFileExists{microtype.sty}{% use microtype if available
  \usepackage[]{microtype}
  \UseMicrotypeSet[protrusion]{basicmath} % disable protrusion for tt fonts
}{}
\makeatletter
\@ifundefined{KOMAClassName}{% if non-KOMA class
  \IfFileExists{parskip.sty}{%
    \usepackage{parskip}
  }{% else
    \setlength{\parindent}{0pt}
    \setlength{\parskip}{6pt plus 2pt minus 1pt}}
}{% if KOMA class
  \KOMAoptions{parskip=half}}
\makeatother
\usepackage{xcolor}
\usepackage[top=0.78in,bottom=0.78in,left=1.22in,right=1.22in]{geometry}
\setlength{\emergencystretch}{3em} % prevent overfull lines
\providecommand{\tightlist}{%
  \setlength{\itemsep}{0pt}\setlength{\parskip}{0pt}}
\setcounter{secnumdepth}{-\maxdimen} % remove section numbering
\setlength{\emergencystretch}{2em}
\setlength{\parskip}{0.55em}
\ifLuaTeX
  \usepackage{selnolig}  % disable illegal ligatures
\fi
\IfFileExists{bookmark.sty}{\usepackage{bookmark}}{\usepackage{hyperref}}
\IfFileExists{xurl.sty}{\usepackage{xurl}}{} % add URL line breaks if available
\urlstyle{same}
\hypersetup{
  pdftitle={From Imagetext to Worldtext},
  pdfauthor={Watson Hartsoe},
  colorlinks=true,
  linkcolor={MidnightBlue},
  filecolor={Maroon},
  citecolor={Blue},
  urlcolor={MidnightBlue},
  pdfcreator={LaTeX via pandoc}}

\title{From Imagetext to Worldtext}
\usepackage{etoolbox}
\makeatletter
\providecommand{\subtitle}[1]{% add subtitle to \maketitle
  \apptocmd{\@title}{\par {\large #1 \par}}{}{}
}
\makeatother
\subtitle{Operative Description and the Persistence of Consequences}
\author{Watson Hartsoe}
\date{Working paper · November 2026}

\begin{document}
\maketitle

\setstretch{1.07}
\textbf{Abstract.} A description can make an image, place an object, or
alter a rule. These operations make different kinds of claims on what
happens next. I call \emph{worldtext} a generative symbolic field in
which words, images, objects, and actions can express, query, and revise
a shared world. Its distinguishing condition is consequential re-entry.
A relation established in one encounter remains available to matter in
another, even when the medium or point of view changes. I develop this
claim through three instruments from my design practice. In the
Builders' Game, a delivery count certifies a threshold that the test
never checks as a doorway. In \emph{House of Dust}, moonlight sets
Penelope and her loom to unraveling, although the animation resets its
phase. In \emph{Shield Grid}, a substring turns a question about
authority into a visual invocation. These cases distinguish execution,
situated adequacy, and persistence. I propose \emph{thick prompting} as
a method for tracing the context, operators, and judgments that let a
description acquire consequences beyond its first output.

\textbf{Keywords} · worldtext, operative description, operative
ekphrasis, thick prompting, generative media, research through design

\hypertarget{the-doorway-that-was-already-complete}{%
\section{1. The doorway that was already
complete}\label{the-doorway-that-was-already-complete}}

The Builders' Game calls one of its challenges THE THRESHOLD. The player
is asked for two pillars and a slab. After those pieces are delivered,
the inspected completion predicate has all it needs. It counts the
deliveries. It does not ask whether the pillars stand apart, whether the
slab spans them, or whether anyone can pass between them. The word
\emph{threshold} suggests an opening that affords passage. What the
program certifies is an inventory. The difference becomes visible
because the apparatus has made one reading of the description executable
(Hartsoe 2026, dossier 1).

This is a useful failure for an account of generative media. The words
have done real work. The pieces are available to be selected and moved.
A partial translation has acquired the authority to announce a completed
task. If the player can return to that arrangement, point to the pieces,
and change what they afford, the initial output becomes part of an
ongoing situation. The capacity to produce the pieces and the capacity
to continue working in the same situation belong to different parts of
the apparatus.

I call that ongoing situation a \emph{worldtext} when it takes the form
of a generative symbolic field that can be described, pictured,
inferred, continued, and re-entered. The field includes the relations
from which its expressions are made and the operations through which
those relations may change. A scene description and a view of the scene
can be co-expressions of a world when a consequential relation can be
recovered across them. Their agreement is an achievement of the
apparatus. A system may achieve it through a runtime, a reader, an
editor, a performer, or a combination of them.

The paper's wager is that \textbf{a description becomes world-forming
when some consequence survives the occasion that produced it and becomes
available to a later encounter}. Survival may mean retaining a fact,
preserving an explicit choice among alternatives, carrying forward a
changed relationship, or keeping a trace of why a revision was made. The
requirement is selective. A world can keep a broken promise while
leaving the color of the door open. That selectivity makes a world both
coherent and unfinished, and gives a participant something specific to
contest.

\hypertarget{from-verbalvisual-relations-to-consequential-re-entry}{%
\section{2. From verbal--visual relations to consequential
re-entry}\label{from-verbalvisual-relations-to-consequential-re-entry}}

The point of departure is \emph{Operation / Describe}, my research
project with Jay David Bolter on the passage from imagetext to Worldtext
(Hartsoe and Bolter n.d.). The account developed here focuses on a
consequence of that passage for design.

Mitchell's \emph{Picture Theory} situates ekphrasis within a larger
inquiry into the relations between the visible and the readable. Its
concept of imagetext offers a way to consider verbal and visual forms in
their conjunction (Mitchell 1994). Bajohr's \emph{operative ekphrasis}
brings a further relation into view. In text-to-image generation, verbal
input participates in producing the image. A caption and a resulting
image can be related through an operation whose material conditions
include encoders, learned representations, and a generation pipeline
(Bajohr 2024a). The joint availability of text and image representations
can enable that operation while the system still assigns them different
roles. In DALL·E 2, a prior predicts an image embedding from text
conditioning and a decoder uses the image embedding to produce pixels
(Ramesh et al.~2022).

Worldtext shifts the question from how a word produces a view to how a
relation established through one expression can bear on another. An
image may show a broken bridge. A later route planner, character, or
editor can treat the break as a fact that narrows the available actions.
Another view may omit the bridge entirely while still belonging to the
same world, provided the break can matter when the route comes into
question. Image fidelity and world continuity therefore answer different
questions. The first asks how an expression presents something. The
second asks what can be carried from it and with what authority.

This extension also draws on the tradition of texts that require action
from their readers. Aarseth's cybertext places the operations that
produce a textual sequence inside the analysis of the work (Aarseth
1997). Navigating a link, executing a score, and editing a scene can
each yield a different realized sequence. Worldtext takes up the
narrower question of what can be carried forward as a commitment when a
sequence becomes material for further production. An editable branch can
retain a decision. An image generator may depend on a reference that
disappears when the resulting image is exported. Interactivity can
coexist with fragile continuity, while a comparatively spare interface
can expose a durable account of what was changed.

Bajohr's \emph{(Berlin, Miami)} makes that division hard to miss. In his
Venice talk he suggests that recurring people, places, and things can
give a reader a world whose causal connections are fragile. His earlier
account calls attention to surface narration, where the form of
storytelling invites readers to infer connections across discontinuities
(Bajohr 2024b, 2025). In his production note he describes reading that
the Jawling had been pregnant, then typing ``Jawling pregnancies'' to
elicit an account of its reproductive cycle (Bajohr 2024c). The two
words were sufficient for the author because he had already identified a
character and a prior event. His reading then became part of the
conditions for the next generation. Here a world relation travels
through a person and back into the system. The published novel alone
cannot reveal how often its apparent continuity depended on that
intervention.

Two computational uses of \emph{world model} help locate the claim. Ha
and Schmidhuber model future observations conditional on an agent's
actions and test transfer from a learned simulator to the original
environment (2018). Wong and colleagues, including Tenenbaum, use a
language model to translate utterances into expressions in a
probabilistic programming language. The resulting program supports
inference and can be extended with new concepts or rules (2023). These
are distinct ways of assigning consequences to an action or utterance.
They help specify what evidence a project should produce when it says a
world remembers, supports intervention, or can be taught a rule.
Worldtext also encompasses the relations that remain with readers,
performers, and editors. A fictional world can be operative in practice
even where its commitments have no unified internal representation.

\hypertarget{locating-the-operator}{%
\section{3. Locating the operator}\label{locating-the-operator}}

The effective unit of operative description is the \emph{context stack}.
It comprises the words, the selected referents, the state available at
the time, the interpretive rules or model, and the validation or
editorial decisions that determine what takes effect. Two people can
utter the same sentence while addressing different entities. The same
person can address the same entity through two interfaces and trigger
different operations. In the inspected Builders' Game, typed ``move
pillar'' reaches the noun-matching delivery loop before the typed move
handler, while ``put that there'' can reach a move operation when the
interface has bound the two referents. The voice route can use an
external model's tool selection. The project thus contains multiple
instruction grammars behind one activity (Hartsoe 2026, dossier 1).

Sol LeWitt's \emph{Wall Drawing 51} gives a precise example of a sparse
description with a thick situation. Its instruction connects the
architectural points of a site by straight lines. A door frame, a
column, or a fire alarm becomes available as a point because it belongs
to the installation's architecture and because a draftsperson can
identify and mark it. The two simultaneous 1970 Turin installations
yielded distinct realizations. The variation belongs to the work's rule,
which calls on what the site supplies (MASS MoCA n.d.). Michael Newman's
reading of Mallarmé's \emph{Le Livre} puts a complementary emphasis on
discretion. In the account accompanying his 2026 Venice talk, an
operator uses chance to place pages in a cabinet under a numerical
constraint. The performer selects an arrangement within a space of
possibilities that someone has already designed (Newman 2026).

These examples distinguish three functions often compressed into the
word \emph{operator}. Someone or something fixes the space of
permissible outcomes. Someone or something chooses a particular
realization. Someone or something decides whether to accept its
consequences. In a generative workflow those functions can be
distributed across an author, a model, a renderer, a runtime, a
platform, and an editor. In LeWitt the physical site also supplies facts
absent from the score. The practical question is which parts of the
context the participant can see and change. The historical question,
posed more broadly by AIMODELS, is how an earlier organization of labor
and measurement comes to be represented as a machine operation (ERC
AIMODELS n.d.).

\hypertarget{when-moonlight-tells-penelope-what-to-do}{%
\section{4. When moonlight tells Penelope what to
do}\label{when-moonlight-tells-penelope-what-to-do}}

Alison Knowles and James Tenney's 1967 \emph{House of Dust} draws
quatrains from four lists, with slots for material, place, light, and
inhabitants. A selected quatrain became the score for a public
sculpture. The poem's circulation between computation and built space
depended on further choices and work (BAMPFA 2022). The Odyssey-based
\emph{House of Dust} in my practice puts that passage between forms
under inspection. Its initial score reads:

\begin{quote}
A HOUSE OF WOVEN THREAD\\
AT THE SWINEHERD'S HUT\\
USING DAYLIGHT\\
INHABITED BY PENELOPE AND A RETURNING STRANGER
\end{quote}

In the implementation, ``a returning stranger'' resolves to the authored
Odysseus asset. This association has a precise provenance. The lookup
carries a reading of the story into the software. It also narrows the
possibilities: a stranger in this score has an assigned identity that
the reader of the quatrain could otherwise leave open. An enumerated
check recorded in the experiment dossier found 500 possible scores, each
of which produced a recipe with three objects and two inhabitants and
survived a round trip through the score encoding. That establishes the
bounded score module's behavior. Loading the required drawings is a
separate operation. The drawing program checks them before committing a
new poem and world together (Hartsoe 2026, dossiers 2--3).

Moonlight is where description begins to act as a rule. In the inspected
source, with the relevant score conditions in place, substituting
moonlight for daylight switches the loom to an unraveling track and
selects Penelope's unweaving pose. The illumination line thus changes
what a person does and what happens to an object. Yet the woven fraction
follows a repeating 16-second phase. A later loop can show the same
action without carrying a material history of how much cloth remains. A
scene can perform a consequence convincingly while the next encounter
inherits only the conditions for performing it again.

This difference suggests a demanding test for Worldtext. Ask the same
question twice, once after the action and once after returning from a
change of view, time, or medium. If the work invites the participant to
believe Penelope has undone a night's labor, there must be some carrier
of that fact when the amount of cloth becomes relevant. An authored loop
can satisfy a different artistic purpose. A production that chooses to
preserve uncertainty might retain two possible amounts or leave the
cloth unmeasured. The test is whether the work makes clear what kind of
continuity the participant is entitled to use.

\hypertarget{when-a-reading-becomes-a-rule}{%
\section{5. When a reading becomes a
rule}\label{when-a-reading-becomes-a-rule}}

\emph{Shield Grid} translates segmented language into visual structure.
In a bounded execution of its extracted classifier, ``Can we decide?''
was assigned QUESTION, while ``Who can decide?'' was assigned
INVOCATION. Both were assigned REPAIR on the world-operation axis. The
question mark initially assigns QUESTION. A later invocation condition
matches the characters ``o'' inside ``Who'' and overwrites that
assignment. Downstream, QUESTION maps to FRACTURE and INVOCATION to
RADIATE (Hartsoe 2026, dossier 5). One substring changes the visual
treatment even though both sentences are questions.

The difference matters because ``Who can decide?'' asks about authority.
The system gives the phrase a different visual force, but the difference
was produced by a lexical match and the ordering of branches. The same
result could be defended as a deliberate reading of the phrase in a
different work. Here the trace identifies what happened in this one. It
also identifies a point for revision. The designer can revise the
matching rule, change how conflicting assignments are handled, or use a
cell override to make a particular reading explicit. Those are different
interventions with different scopes.

The case identifies three judgments that a fluent output can compress.
\textbf{Execution} asks which operation occurred. \textbf{Adequacy} asks
whether the operation fulfills the situated task or supports the
interpretation claimed for it. \textbf{Persistence} asks which of its
effects continue to constrain later encounters. The threshold tally
executes and can be retained, while its adequacy as a doorway test
remains open. Penelope's unweaving is legible within a scene, while the
amount undone need not persist. A glyph assignment may propagate
reliably even when the interpretation behind it is fragile. An account
of Worldtext has to leave room for disagreement at each of these points.

\hypertarget{thick-prompting-as-a-method}{%
\section{6. Thick prompting as a
method}\label{thick-prompting-as-a-method}}

Geertz's use of thick description, following Ryle's wink example,
distinguishes bodily movement from the socially situated action it may
perform (Geertz 1973). A longer account of the eyelid would leave the
relevant difference unexplained. Thick prompting borrows that attention
to consequential distinctions. It studies the relation between an
utterance and what a practice makes possible with it. A two-word phrase
can be thick when it invokes a character with a recoverable history. A
page of elaborate description can be thin as evidence if the conditions
that made its output possible are unavailable.

The unit of evidence is a documented prompt--output sequence. Preserve
the wording, the addressed entities, the initial state, the route
through the system, any supplied references or rules, the resulting
artifact, and the decision that accepted or rejected it. Then vary one
consequential condition. Keep the parts constant and change their
placement. Keep the score constant and replace daylight with moonlight.
Keep the punctuation and change the phrasing. Follow the result through
a second encounter and record which effect survives, who can revise it,
and how that revision is attributed. When a model is stochastic, record
the model and sampling conditions when available and retain a bounded
set of alternatives. A comparison supports claims about the observed
conditions and trajectories.

The Association Field in Cineosis offers a different kind of sequence.
An editor walks a path of clips and uses LAY to seat them into a spoken
line. The inspected implementation uses line boundaries supplied by its
kernel and can carry remaining steps forward when the current line
cannot accommodate the full path (Hartsoe 2026, dossier 4). The same
words and source clips can therefore support different cuts. What a
later editor needs to recover is the chosen path, the resulting
intervals, and the distinction between a candidate clip and a committed
edit. A rendered montage can carry the composition to an audience while
leaving these decisions inaccessible to the next maker. The ability to
re-enter a world is consequently also a question of continuing
authorship.

One useful stopping rule is to end a comparison once there is enough
evidence to distinguish the proposed explanations. In Shield Grid, the
extracted rule and its contrasting inputs identify a mechanism. More
attractive renderings would add little to that finding. A claim about
what viewers understand would require a different comparison. In a
creative workflow, the aim may instead be to preserve several materially
different readings so that the next maker can choose between them. This
is why a record of a rejected interpretation can matter as much as a
prompt that worked. It gives the world a history that someone can
reopen.

\hypertarget{limits-and-implications}{%
\section{7. Limits and implications}\label{limits-and-implications}}

The cases support distinctions among specific operations. They are based
on source inspection and bounded executions recorded in the project
dossiers. The doorway comparison remains a proposed run, the loom
evidence concerns a repeating phase, and the Shield Grid check exercised
an extracted function. These materials do not establish how participants
understand the interfaces or how a particular generative model preserves
state over long runs. A production account, such as Villeneuve's
description of puppeteers creating shadows for the performers in
\emph{Arrival}, offers an adjacent case of a deliberately arranged cue.
It shows how a world can be made available to someone through timing and
presence before it is available to an audience as a finished image. The
contribution of that cue to the finished performance remains a separate
inference (Villeneuve and Beguš 2026, 15:25--17:30).

A strong objection is that insisting on persistence would exclude
literary and artistic worlds whose discontinuities are essential to
them. The objection helps specify what is meant by a consequence. A
contradiction can persist as a contradiction. Two versions of an event
can remain separately available and shape different encounters. A
participant may be entitled to inspect which version a scene uses
without having the power to resolve the dispute. Worldtext can be
incomplete and disputable because what it must preserve is the
availability of a relation for further interpretation or action. An
interface that silently overwrites the losing version makes a different
artistic and epistemic choice from one that allows both to be revisited.

The term also does not require treating a model's latent space as a
complete fictional universe. A latent representation can contribute
productive possibilities. A canon can establish selected commitments. An
artifact can express some of them, and a history can record how they
changed. Those are different functions, even when one system performs
several of them. Asking where each function lives makes it possible to
compare worlds assembled through human practice, authored rules, learned
models, and combinations of the three.

\hypertarget{conclusion}{%
\section{Conclusion}\label{conclusion}}

An operative description can give a world something to keep. The
Builders' Game shows how narrowly a machine may certify what has
happened. \emph{House of Dust} shows how much a line of poetry can set
in motion and how little of that motion a loop must remember.
\emph{Shield Grid} shows how a small reading decision can acquire
visible force and why its provenance matters. In a worldtext, the next
description arrives in a situation changed by the last one. Its new
referents, possible actions, and rights of revision are the substance of
the world.

\clearpage

\hypertarget{references}{%
\section{References}\label{references}}

\small

Aarseth, Espen J. 1997.
\href{https://www.press.jhu.edu/books/title/1398/cybertext}{\emph{Cybertext:
Perspectives on Ergodic Literature}}. Johns Hopkins University Press.

Bajohr, Hannes. 2024a. ``Operative Ekphrasis: The Collapse of the
Text/Image Distinction in Multimodal AI.'' \emph{Word \& Image} 40 (2),
77--90.
\href{https://doi.org/10.1080/02666286.2024.2330335}{doi:10.1080/02666286.2024.2330335}.

Bajohr, Hannes. 2024b.
``\href{https://transit.berkeley.edu/2024/cohesion-without-coherence-artificial-intelligence-and-narrative-form/}{Cohesion
Without Coherence: Artificial Intelligence and Narrative Form}.''
Translated by Kayla Rose van Kooten. \emph{TRANSIT} 14 (2).

Bajohr, Hannes. 2024c. ``(Berlin, Miami).'' Excerpts and production
note.
\href{https://www.ensemblepark.com/issues/ensemble_park_1.pdf}{\emph{Ensemble
Park} 1}.

Bajohr, Hannes. 2025.
``\href{https://www.youtube.com/watch?v=hmchpa_Ek20}{The Latent Space of
Meaning and the Novel: World Models in AI and Literature}.'' Chat Token
Vector, Ca' Foscari University of Venice, June 12.

BAMPFA. 2022.
\href{https://bampfa.org/sites/default/files/AlisonKnowles_LargePrint_DIGITAL_v3.pdf}{\emph{By
Alison Knowles: A Retrospective (1960--2022)}}. Exhibition text.

ERC AIMODELS. n.d.
``\href{https://pric.unive.it/projects/ai-models/home}{The Culture of
Algorithmic Models: Advancing the Historical Epistemology of Artificial
Intelligence}.'' Ca' Foscari University of Venice. Project description.

Geertz, Clifford. 1973. ``Thick Description: Toward an Interpretive
Theory of Culture.'' In \emph{The Interpretation of Cultures}. Basic
Books.

Ha, David, and Jürgen Schmidhuber. 2018.
``\href{https://worldmodels.github.io/}{World Models}.''
arXiv:1803.10122.

Hartsoe, Watson. 2026.
``\href{https://chatgpt.com/library?file_id=libfile_6a9f9eac7dfc8191b18354ff8742c3f5\&account_id=personal}{Grounded
Experiment Dossiers: From Wording to Consequence}.'' Research notes,
October 15. Primary artifacts:
\href{https://hartswf0.github.io/tractor-dce-gyo/builders-game.html}{Builders'
Game},
\href{https://hartswf0.github.io/odyssey-halfworld/house-of-dust.html}{House
of Dust}, and
\href{https://hartswf0.github.io/text-ure/shield\%2052/1.html}{Shield
Grid}.

Hartsoe, Watson, and Jay David Bolter. n.d.
\href{https://hartswf0.github.io/operation-describe/}{\emph{Operation /
Describe: From Imagetext to Worldtext}}. Research project.

MASS MoCA. n.d. ``\href{https://massmoca.org/event/walldrawing51/}{Wall
Drawing 51}.'' Exhibition entry for Sol LeWitt's June 1970 work.

Mitchell, W. J. T. 1994.
\href{https://press.uchicago.edu/ucp/books/book/chicago/P/bo3683962.html}{\emph{Picture
Theory: Essays on Verbal and Visual Representation}}. University of
Chicago Press.

Newman, Michael. 2026.
``\href{https://www.youtube.com/watch?v=ce3rNpXdu30}{Art as `Operation':
From Lovelace and Mallarmé to Adversarial Poetry}.'' Recording
description. Adversarial Poetry symposium, Ca' Foscari University of
Venice, October 1.

Ramesh, Aditya, Prafulla Dhariwal, Alex Nichol, Casey Chu, and Mark
Chen. 2022. ``\href{https://arxiv.org/abs/2204.06125}{Hierarchical
Text-Conditional Image Generation with CLIP Latents}.''
arXiv:2204.06125.

Villeneuve, Denis, and Nina Beguš. 2026.
``\href{https://www.youtube.com/watch?v=7QoS5uK52DY}{Arrival: A
Conversation with Denis Villeneuve and Dr.~Nina Beguš}.'' Academy Museum
of Motion Pictures, November 2.

Wong, Lionel, Gabriel Grand, Alexander K. Lew, Noah D. Goodman, Vikash
K. Mansinghka, Jacob Andreas, and Joshua B. Tenenbaum. 2023.
``\href{https://arxiv.org/abs/2306.12672}{From Word Models to World
Models: Translating from Natural Language to the Probabilistic Language
of Thought}.'' arXiv:2306.12672.

\end{document}
